\documentclass{article}
\usepackage{xintexpr}
\usepackage{booktabs}
\usepackage{siunitx}
\usepackage{geometry}
\geometry{margin=1in}

% Con 8 decimales para mejor visualización
\newcommand{\miintegral}[2]{%
	\xintthefloatexpr[8]
	add(exp(-(i*#1/#2)^2), i=1..(#2-1)) * (#1/#2) + 
	(exp(0) + exp(-(#1)^2)) * (#1/(2*#2))
	\relax
}
\newcommand{\errorabs}[3]{%
	\xintthefloatexpr[10]
	abs(#1 - \miintegral{#2}{#3})
	\relax
}
\begin{document}
	
\section{Regla del trapecio con el paquete \texttt{xintexpr}}	
Definimos la función $f(x) = \displaystyle\int_0^x e^{-t^2} dt$, usando la regla del trapecio	$n=20$ pasos de integración\\

\bigskip	
\textbf{Convergencia de la regla del trapecio}\\
\begin{tabular}{cccccc}
	\toprule
	$x$ & $n=6$ & $n=8$ & $n=10$ & $n=12$ & $n=14$ \\ \midrule
	0.5 & \miintegral{0.5}{6} & \miintegral{0.5}{8} & \miintegral{0.5}{10} & \miintegral{0.5}{12} & \miintegral{0.5}{14} \\
	1.0 & \miintegral{1.0}{6} & \miintegral{1.0}{8} & \miintegral{1.0}{10} & \miintegral{1.0}{12} & \miintegral{1.0}{14} \\
	1.5 & \miintegral{1.5}{6} & \miintegral{1.5}{8} & \miintegral{1.5}{10} & \miintegral{1.5}{12} & \miintegral{1.5}{14} \\ \bottomrule
\end{tabular}

\bigskip
\textbf{Errores aproximados respecto al valor de referencia $f(0.5)=0.46128100641279246$}\\

\begin{tabular}{cS[table-format=1.2e-2]}
	\toprule
	$n$ & \multicolumn{1}{c}{Error absoluto ($x=0.5$)} \\ \midrule
	6 & \errorabs{0.46128100641279246}{0.5}{6} \\
	8 & \errorabs{0.46128100641279246}{0.5}{8} \\
	10 & \errorabs{0.46128100641279246}{0.5}{10} \\
	12 & \errorabs{0.46128100641279246}{0.5}{12} \\
	14 & \errorabs{0.46128100641279246}{0.5}{14} \\ \bottomrule
\end{tabular}
	
\bigskip	
\textbf{Resumen de la sintaxis usada con el paquete \texttt{xintexpr}}
	
	\begin{enumerate}
		\item \verb+\xintthefloatexpr+: inicia la evaluación de una expresión en punto flotante
		\item \verb+\relax+: termina la expresión (alternativa: \verb+\xintendexpr+)
		\item \verb+add(expr, variable = inicio .. fin)+: suma iterativa
		\item \verb+exp(x)+: función exponencial $e^x$
		\item \verb+i*#1/20+: operaciones aritméticas básicas
	\end{enumerate}
	
\end{document}

\documentclass{article}
\usepackage{xintexpr}
\usepackage{booktabs}
\usepackage{geometry}
\geometry{margin=1in}

% Con 8 decimales para mejor visualización
\newcommand{\miintegral}[2]{%
	\xintthefloatexpr[8]
	add(exp(-(i*#1/#2)^2), i=1..(#2-1)) * (#1/#2) + 
	(exp(0) + exp(-(#1)^2)) * (#1/(2*#2))
	\relax
}

\begin{document}
	
	\section{Convergencia de la regla del trapecio}
	
	\begin{tabular}{cccccc}
		\toprule
		$x$ & $n=6$ & $n=8$ & $n=10$ & $n=12$ & $n=14$ \\ \midrule
		0.5 & \miintegral{0.5}{6} & \miintegral{0.5}{8} & \miintegral{0.5}{10} & \miintegral{0.5}{12} & \miintegral{0.5}{14} \\
		1.0 & \miintegral{1.0}{6} & \miintegral{1.0}{8} & \miintegral{1.0}{10} & \miintegral{1.0}{12} & \miintegral{1.0}{14} \\
		1.5 & \miintegral{1.5}{6} & \miintegral{1.5}{8} & \miintegral{1.5}{10} & \miintegral{1.5}{12} & \miintegral{1.5}{14} \\ \bottomrule
	\end{tabular}
	
	\subsection*{Errores aproximados respecto al valor de referencia ($x=0.5$, ref = 0.5204998778)}
	
\begin{tabular}{cS[table-format=1.2e-2]}
	\toprule
	$n$ & \multicolumn{1}{c}{Error absoluto ($x=0.5$)} \\ \midrule
	6 & \errorabs{0.46128100641279246}{0.5}{6} \\
	8 & \errorabs{0.46128100641279246}{0.5}{8} \\
	10 & \errorabs{0.46128100641279246}{0.5}{10} \\
	12 & \errorabs{0.46128100641279246}{0.5}{12} \\
	14 & \errorabs{0.46128100641279246}{0.5}{14} \\ \bottomrule
\end{tabular}
	
\end{document}